\cotefile{157-3}
\scanpage{34}
\mspage{4}

$W_{\mathbb D}$. Les conditions L1 à L5 sont \unsure{énoncées}
dans \unsure{PS|P5} § 64 (page 164).

\uline{L1} Saturation de $W_{\mathbb D}$ O.K.

\deleted{\unsure{en}} \deleted{\unsure{f}}Cela résulte \inserted{\unsure{formellt}} \deleted{du fait} que l'on \deleted{\unsure{fait}} a
un foncteur
\[
\mathrm{Cat}\xrightarrow{\varphi} C\qquad (=\mathrm{Cat}^{\circ})
\]
tel que \unsure{pour} $u\in\mathrm{Fl}(\mathrm{Cat})$, $u\in W_{\mathbb D}\Longleftrightarrow\unsure{\varphi(u)|\varphi u}$ iso.
(Prendre pour $\varphi(X)$ la \unsure{sous-catégorie} de
$\deleted{\unsure{A}}\mathbb D(X)$ formée des \deleted{\unsure{faisceaux}} \inserted{$\mathbb D$-}coeff. \uline{constants} sur $X$.)

\uline{L2} ``$\mathbb{\Delta}_1$ est $W_{\deleted{\unsure{\mathbb D}}}$-asphérique sur $e$'', i.e. pour tt
$X$ dans $\mathrm{Cat}$, $X\times\mathbb{\Delta}_1\longrightarrow X$ est dans $W$.

Cela signifie aussi que $\mathbb{\Delta}_1\longrightarrow e$ est une $\mathbb D$-fibration
triviale, ce qui est bien le cas.

\uline{L3}, \deleted{(\unsure{renforce} L2) :} Si $X\in\mathrm{Ob}\,\mathrm{Cat}$ a un
objet final, alors $X\longrightarrow e$ est \deleted{$W$-asphérique} \inserted{dans $W$}

% slanted
\margin{En fait $L2\unsure{\Longrightarrow|=}L3$}
\deleted{i.e.} En fait, nous exigeons pour les dérivateurs
que ceci soit vrai \deleted{\unsure{pour}}\inserted{en supposant seulement} $X$ asphérique
(au sens absolu), \deleted{\unsure{seulement},} et on en
conclut \inserted{même} que $X\longrightarrow e$ est $W$-asphérique
i.e. $X\times Y\xrightarrow{\mathrm{pr}_2}Y$ est dans $W$ pour tt $Y$
dans $\mathrm{Cat}$.

\uline{L4} Si $f:X\longrightarrow Y$ est tel que $X_{/y}\longrightarrow Y_{/y}$
sont dans $W$, alors $f$ est dans $W$.
C'est \unsure{OK|OV}.

% archive sheet number: 33
